\documentclass[11pt,a4paper]{article}
\usepackage[utf8]{inputenc}
\usepackage[margin=2cm]{geometry}
\usepackage{amsmath,amssymb,amsfonts}
\usepackage{enumitem}
\usepackage{tcolorbox}

\tcbset{
    colback=gray!5,
    colframe=black!75,
    boxrule=0.6pt,
    arc=1.5mm,
    left=2.5mm,
    right=2.5mm,
    top=1.5mm,
    bottom=1.5mm
}

\setlength{\parindent}{0pt}
\setlength{\parskip}{0.4em}

\title{\vspace{-1.5cm}\textbf{Solusi Tugas Latihan PTSB}\\\large Pengantar Teori Suku Bunga --- 7 Oktober 2026\vspace{-0.8cm}}
\author{}
\date{}

\begin{document}

\maketitle

\vspace{-0.5em}
\hrule
\vspace{0.8em}

% =============================================================
% NOMOR 1
% =============================================================
\begin{tcolorbox}
\textbf{1.} Bila diketahui $i^{(m)} = a$ dan $d^{(m)} = b$, tentukan $m$ dalam $a$ dan $b$.
\end{tcolorbox}

\textbf{Solusi:}

Berdasarkan definisi suku bunga nominal dan tingkat diskonto nominal:
\begin{align*}
    1 + i &= \left(1 + \frac{a}{m}\right)^m \\
    1 - d &= \left(1 - \frac{b}{m}\right)^m
\end{align*}

Dari hubungan $(1 + i)(1 - d) = 1$:
\begin{align*}
    \left(1 + \frac{a}{m}\right)^m \left(1 - \frac{b}{m}\right)^m &= 1 \\[0.3em]
    \left(1 + \frac{a}{m}\right)\left(1 - \frac{b}{m}\right) &= 1 \\[0.3em]
    1 - \frac{b}{m} + \frac{a}{m} - \frac{ab}{m^2} &= 1 \\[0.3em]
    \frac{a - b}{m} - \frac{ab}{m^2} &= 0
\end{align*}

Kalikan kedua ruas dengan $m^2$:
\begin{align*}
    (a - b)m - ab &= 0 \\[0.3em]
    (a - b)m &= ab \\[0.3em]
    \mathbf{m} &= \mathbf{\frac{ab}{a - b}}
\end{align*}

\vspace{0.5em}

% =============================================================
% NOMOR 3
% =============================================================
\begin{tcolorbox}
\textbf{3.} Irfan memiliki modal sebesar 1 juta Rupiah. Ia memiliki 2 opsi untuk melakukan investasi yaitu ikut dalam Projek A atau Projek B. Kedua projek tersebut memerlukan modal yang sama yaitu 1 juta Rupiah. Jika berinvestasi pada Projek A, Irfan akan mendapatkan 1.5 juta Rupiah di akhir tahun pertama. Sedangkan projek B akan memberikan cashflow sebesar 750 ribu Rupiah di akhir tahun pertama dan 1 juta Rupiah di akhir tahun kedua. Tentukan IRR Projek A dan Projek B, kemudian beri keputusan sebaiknya Irfan melakukan investasi pada Projek A atau B!
\end{tcolorbox}

\textbf{Solusi:}

\textbf{A. IRR Projek A ($r_A$):}
\begin{itemize}[leftmargin=1.5em, topsep=0pt, itemsep=0pt]
    \item $C_0 = -1\text{ juta}$, $C_1 = +1.5\text{ juta}$.
\end{itemize}
\begin{equation*}
    NPV_A(r_A) = -1 + \frac{1.5}{1 + r_A} = 0 \implies 1 + r_A = 1.5 \implies \mathbf{r_A = 50\%}
\end{equation*}

\textbf{B. IRR Projek B ($r_B$):}
\begin{itemize}[leftmargin=1.5em, topsep=0pt, itemsep=0pt]
    \item $C_0 = -1\text{ juta}$, $C_1 = +0.75\text{ juta}$, $C_2 = +1\text{ juta}$.
\end{itemize}
\begin{equation*}
    NPV_B(r_B) = -1 + \frac{0.75}{1 + r_B} + \frac{1}{(1 + r_B)^2} = 0
\end{equation*}
Misalkan $v = \frac{1}{1 + r_B}$:
\begin{align*}
    -1 + 0.75v + v^2 &= 0 \iff 4v^2 + 3v - 4 = 0 \\[0.3em]
    v_{1,2} &= \frac{-3 \pm \sqrt{3^2 - 4(4)(-4)}}{2(4)} = \frac{-3 \pm \sqrt{73}}{8}
\end{align*}
Dua kemungkinan nilai $v$:
\begin{itemize}[leftmargin=1.5em, topsep=0pt, itemsep=0pt]
    \item $v_1 = \dfrac{-3 + \sqrt{73}}{8} \approx \mathbf{0.6930005}$ (valid, $v > 0$) $\implies 1 + r_B = \dfrac{1}{v_1} = \dfrac{\sqrt{73} + 3}{8} \approx 1.4430005 \implies \mathbf{r_B = 44.30\%}$.
    \item $v_2 = \dfrac{-3 - \sqrt{73}}{8} \approx \mathbf{-1.4430005}$ (tidak valid secara finansial karena $v < 0$, ditolak).
\end{itemize}

\textbf{C. Keputusan:}
\begin{equation*}
    IRR_A \, (50\%) > IRR_B \, (44.30\%)
\end{equation*}
Karena modal awal kedua proyek bernilai sama dan $IRR_A > IRR_B$, maka sebaiknya Irfan memilih \textbf{Projek A}.

\vspace{0.5em}

% =============================================================
% NOMOR 5
% =============================================================
\begin{tcolorbox}
\textbf{5.} Tentukan present value dari 10-year annuity due dari pembayaran \$1 per tahun, jika suku bunga efektif 4\% untuk 5 tahun pertama dan 5\% untuk 5 tahun berikutnya!
\end{tcolorbox}

\textbf{Solusi:}

Pembayaran $\$1$ terjadi pada $t = 0, 1, 2, \dots, 9$:
\begin{itemize}[leftmargin=1.5em, topsep=0pt, itemsep=0pt]
    \item 5 tahun pertama ($t = 0, 1, 2, 3, 4$): $i_1 = 4\%$, $v_1 = (1.04)^{-1}$.
    \item 5 tahun kedua ($t = 5, 6, 7, 8, 9$): $i_2 = 5\%$, $v_2 = (1.05)^{-1}$.
\end{itemize}

Persamaan Nilai Sekarang ($PV$) pada $t = 0$:
\begin{align*}
    PV &= \sum_{t=0}^4 v_1^t + v_1^5 \sum_{k=0}^4 v_2^k \\
    &= \ddot{a}_{\overline{5}|0.04} + v_1^5 \cdot \ddot{a}_{\overline{5}|0.05}
\end{align*}

Perhitungan nilai komponen:
\begin{align*}
    \ddot{a}_{\overline{5}|0.04} &= (1.04) \left[ \frac{1 - (1.04)^{-5}}{0.04} \right] = (1.04)(4.451822) = 4.629895 \\[0.3em]
    \ddot{a}_{\overline{5}|0.05} &= (1.05) \left[ \frac{1 - (1.05)^{-5}}{0.05} \right] = (1.05)(4.329477) = 4.545951 \\[0.3em]
    v_1^5 \cdot \ddot{a}_{\overline{5}|0.05} &= (1.04)^{-5} \times 4.545951 = (0.821927) \times 4.545951 = 3.736440
\end{align*}

Maka:
\begin{equation*}
    PV = 4.629895 + 3.736440 = \mathbf{8.366335 \approx \$8.37}
\end{equation*}

\end{document}
